% Section II (W1). Every technical sentence ends with its source tag.
% Rival descriptions come from facts-F8 only; the "this work" row of
% Table tab:related comes from facts-live/F1/F2/F3 (see paper/data/make_related.py).
\section{Related Work}\label{sec:related}

Table~\ref{tab:related} compares the closest work with ours, group by group.
% src: facts-F8#23

% Generated by paper/data/make_related.py from paper/notes/facts-F8.md
% (rival rows) and facts-live/F1/F2/F3/F8 (this-work row). Do not edit by hand.
\begin{table}[t]
\centering
\caption{Comparison with related work (legend and notes below the table).}
\label{tab:related}
\footnotesize
\setlength{\tabcolsep}{2.4pt}
\begin{tabular}{@{}lcccccccc@{}}
\toprule
Work & CE & LG & BG & LC & TB & ES & FP & IMP \\
\midrule
Kroll et al.~\cite{kroll2018accountable} & $\checkmark$ & $\checkmark$ & $\circ$ & -- & ? & -- & $\circ$ & $\checkmark$ \\ % src: facts-F8#1
Ryan~\cite{ryan2017accountable} & $\checkmark$ & $\checkmark$ & -- & -- & -- & -- & $\circ$ & -- \\ % src: facts-F8#2
Cheval et al.~\cite{cheval2023transparency} & $\checkmark$ & $\checkmark$ & -- & -- & -- & -- & $\checkmark$ & $\circ$ \\ % src: facts-F8#3
Larch~\cite{dauterman2023larch} & $\checkmark$ & $\checkmark$ & -- & -- & -- & -- & ? & $\checkmark$ \\ % src: facts-F8#4
Li et al.~\cite{li2025accountable} & $\checkmark$ & $\checkmark$ & ? & -- & -- & -- & $\checkmark$ & $\checkmark$ \\ % src: facts-F8#5
\midrule
Brucker, Petritsch~\cite{brucker2009extending} & -- & -- & $\checkmark$ & -- & ? & -- & ? & ? \\ % src: facts-F8#6
Ferreira et al.~\cite{ferreira2006break} & -- & -- & $\checkmark$ & -- & ? & -- & ? & $\checkmark$ \\ % src: facts-F8#7
Marinovic et al.~\cite{marinovic2011rumpole} & -- & ? & $\checkmark$ & ? & ? & -- & ? & ? \\ % src: facts-F8#8
\midrule
Brucker et al.~\cite{brucker2010abe} & $\checkmark$ & ? & $\checkmark$ & ? & ? & -- & ? & ? \\ % src: facts-F8#9
Zhang et al.~\cite{zhang2016password} & $\checkmark$ & -- & $\checkmark$ & ? & ? & -- & ? & ? \\ % src: facts-F8#10
Scafuro~\cite{scafuro2019breakglass} & $\checkmark$ & $\circ$ & $\checkmark$ & -- & -- & -- & $\checkmark$ & -- \\ % src: facts-F8#11
de Oliveira et al.~\cite{deoliveira2019redalert} & $\checkmark$ & -- & $\checkmark$ & $\circ$ & $\checkmark$ & -- & -- & -- \\ % src: facts-F8#12
de Oliveira et al.~\cite{deoliveira2020breakglass} & $\checkmark$ & -- & $\checkmark$ & $\circ$ & $\checkmark$ & -- & $\checkmark$ & $\checkmark$ \\ % src: facts-F8#13
LiBAC~\cite{yang2018libac} & $\checkmark$ & -- & $\checkmark$ & -- & ? & -- & $\checkmark$ & $\checkmark$ \\ % src: facts-F8#14
\midrule
Li et al.~\cite{li2013phr} & $\checkmark$ & -- & $\checkmark$ & $\circ$ & -- & -- & ? & $\checkmark$ \\ % src: facts-F8#15
\midrule
Yang, Ma~\cite{yang2016timing} & $\checkmark$ & -- & -- & $\circ$ & $\checkmark$ & $\checkmark$ & $\checkmark$ & $\circ$ \\ % src: facts-F8#16
\midrule
Muslim, Gr\'egoire~\cite{muslim2026consent} & -- & -- & $\checkmark$ & $\circ$ & ? & -- & $\circ$ & $\checkmark$ \\ % src: facts-F8#17
Romel et al.~\cite{romel2025patient} & $\checkmark$ & -- & $\checkmark$ & $\circ$ & $\checkmark$ & -- & $\circ$ & $\checkmark$ \\ % src: facts-F8#18
\midrule
\textbf{This work} & $\checkmark$ & $\checkmark^{a}$ & $\checkmark$ & $\circ^{b}$ & $\checkmark^{c}$ & $\checkmark$ & -- & $\checkmark$ \\ % src: facts-live#12, facts-F1#11, facts-F1#35, facts-F1#54, facts-F1#77, facts-F1#78, facts-F2#50, facts-F2#59, facts-F3#53, facts-F8#3
\bottomrule
\end{tabular}
\par\vspace{3pt}
{\raggedright CE: keys, not a policy check, decide access; LG: key release
gated on a log entry; BG: break-glass; LC: consent and revocation in one
lifecycle; TB: time-boxed expiry; ES: emergency-scoped or time-limited
search; FP: formal proof or model; IMP: implementation. $\checkmark$ yes,
$\circ$ partial, -- no, ? not determined from what we read. Our row: live
run on \texttt{57a1a8c}.
% src: facts-live#12
$^{a}$Trusted hash chain, no log proofs.
% src: facts-F2#59, facts-F8#3
$^{b}$Revocation is a transition; consent is a signed grant checked at
read time when presented, outside the lifecycle.
% src: facts-F1#11, facts-F1#35, facts-F1#77, facts-F1#78
$^{c}$Server-enforced, not cryptographic.\par}
% src: facts-F1#54
\end{table}


\textit{Accountable and transparent decryption.}
Ryan~\cite{ryan2017accountable} proposes accountable decryption, in which a trusted hardware device decrypts only after the request is entered in a log whose untrusted maintainer gives proofs of presence and extension.
% src: facts-F8#2
Cheval, Moreira and Ryan~\cite{cheval2023transparency} formally verify transparent decryption in ProVerif: trustees hold shares of the decryption key and use them only if the Merkle-tree log proofs verify.
% src: facts-F8#3
Kroll et al.~\cite{kroll2018accountable} mediate every access through an approver and a threshold of decryption authorities, and the other parties do not proceed until an auditor certifies that the access is in its publicly verifiable log.
% src: facts-F8#1
They name hospital emergencies as an application: an unpredictable set of staff may suddenly need a patient's records, and the audit record lets internal reviewers check afterwards that the access was not abused~\cite{kroll2018accountable}.
% src: facts-F8#1
Larch~\cite{dauterman2023larch} applies the log-as-secret-holder pattern to authentication, and Li et al.~\cite{li2025accountable} give formal definitions of accountable decryption and a TEE-assisted construction.
% src: facts-F8#4, facts-F8#5
Transparent decryption is the primitive we instantiate: our logger holds two of four emergency-key shares, checks a signed Emergency Assertion before it commits, and sits inside a revocation and expiry lifecycle.
% src: facts-F8#1, facts-F8#3, facts-F2#34, facts-F2#36, facts-F2#45, facts-F1#47
Our log is weaker than all of these: a trusted hash chain, not publicly verifiable, with no proofs and no formal model.
% src: facts-F8#1, facts-F8#3, facts-F2#59

\textit{Break-glass access control.}
Ferreira et al.~\cite{ferreira2006break} add a break-the-glass policy with non-repudiation of its use to the access-control model of a hospital EMR, while Brucker and Petritsch~\cite{brucker2009extending} and Marinovic et al.~\cite{marinovic2011rumpole} build break-glass into access-control models.
% src: facts-F8#6, facts-F8#7, facts-F8#8
These are policy checks in software; in ours, keys and the threshold gate decide access, and the log entry precedes key release.
% src: facts-F8#6, facts-F8#7, facts-F4#19, facts-F2#51

\textit{Cryptographic break-glass.}
Brucker et al.~\cite{brucker2010abe} add break-glass to ABE-based information sharing, LiBAC~\cite{yang2018libac} lets an emergency override a file's ABE policy in healthcare IoT, and Zhang et al.~\cite{zhang2016password} let a password holder obtain a decryption key from a key generation centre that is held accountable for abuse.
% src: facts-F8#9, facts-F8#14, facts-F8#10
Scafuro~\cite{scafuro2019breakglass} builds encryption breakable exactly once, detectably, from hardware tokens, as a feasibility result.
% src: facts-F8#11
De Oliveira et al.\ limit emergency access to the task at hand~\cite{deoliveira2019redalert}, or to each team's period of care across institutions using signed access tokens and revocation by a trusted master authority and storage provider, without re-encrypting data~\cite{deoliveira2020breakglass}.
% src: facts-F8#12, facts-F8#13
None gates key release on a committed log entry (Scafuro's break is detectable, not log-gated), and our expiry re-keys the record instead of trusting a provider to block access; several, however, have formal analyses we lack.
% src: facts-F8#9, facts-F8#10, facts-F8#11, facts-F8#12, facts-F8#13, facts-F8#14, facts-F1#39

\textit{Encrypted PHR sharing.}
Li et al.~\cite{li2013phr} share personal health records on semi-trusted servers with multi-authority ABE, with multiple security domains, on-demand user and attribute revocation, and break-glass access in emergencies.
% src: facts-F8#15
We keep their patient-centric, multi-authority setting, but release the emergency key only through an audit-logged threshold, and the emergency ends with a rotation.
% src: facts-F8#15, facts-F2#36, facts-F1#55

\textit{Time-limited search.}
Yang and Ma~\cite{yang2016timing} let patients delegate keyword search and decryption rights over their records for a controlled period.
% src: facts-F8#16
This is prior art; we only derive emergency search tokens from the emergency key, so search ends at the expiry rotation, an engineering choice.
% src: facts-F8#16, facts-F3#51, facts-F3#54

\textit{Consent at the policy layer.}
Muslim and Gr\'egoire~\cite{muslim2026consent} validate consent directives before activation in XACML, add bounded break-the-glass access logged for post-hoc review, and argue that consent needs a lifecycle of creation, activation, revocation and expiration; their enforcement is ``logical rather than cryptographic''.
% src: facts-F8#17
We realise revocation and expiration by re-keying, while consent stays a patient-signed grant checked at read time; their ``pre-commit'' commits a policy, not a log entry.
% src: facts-F8#17, facts-F1#47, facts-F1#77, facts-F1#78
Romel et al.~\cite{romel2025patient} keep encrypted records off-chain, with patient-signed, time-bounded permissions on a public chain and keys distributed by public-key wrapping; revocation flips an on-chain flag, and re-keying is a separate, optional workflow.
% src: facts-F8#18
They also support emergency access: two registered emergency physicians co-sign a request that expires after two hours and is emitted as an on-chain event; that grant is logged, but key release goes through a separate, design-only HSM-backed guardian, not gated on the log, whereas ours needs the audit service's shares.
% src: facts-F8#18, facts-F2#36
